\documentclass[11pt]{article}

\usepackage{amsmath,amssymb}
\usepackage{xurl}
\usepackage[hidelinks]{hyperref}
\setlength{\emergencystretch}{2em}

% Shared commands for notes in docs/paper-gaps/.
%
% Two halves.  The link commands below are project identity: OWNER/REPO is the
% placeholder that scripts/bootstrap_project.py replaces with this project's
% GitHub slug and Pages site.  Everything after them is NOTATION, and notation
% belongs to the source paper: the definitions here are an example set, and a
% project replaces them with the macros of its own paper's style file.  The one
% rule that never changes: a command defined here must mean exactly what the
% source paper's macro means, or a note silently says something else.

\newcommand{\Lean}{\textsc{Lean}}
\newcommand{\leanid}[1]{\textsf{#1}}
\newcommand{\ghissue}[1]{\href{https://github.com/OWNER/REPO/issues/#1}{GitHub issue \##1}}
\newcommand{\ghpr}[1]{\href{https://github.com/OWNER/REPO/pull/#1}{PR \##1}}
% Local issue tracker (issues/NNNN-*.md in this repository); no remote URL.
\newcommand{\localissue}[1]{issue \##1 (\path{issues/})}
\newcommand{\doclink}[1]{\href{https://OWNER.github.io/REPO/docs/#1.html}{\path{#1}}}

% --- Notation: replace with the source paper's own macros --------------------
% \F, \N, \C, \R, \tr, \ot, \eps, \E, \bx, \bu, \bv, \by

\newcommand{\F}{\mathbb{F}}
\newcommand{\N}{\mathbb{N}}
\newcommand{\C}{\mathbb{C}}
\newcommand{\R}{\mathbb{R}}
\newcommand{\tr}{\mathrm{tr}}
\newcommand{\eps}{\epsilon}
\newcommand{\ot}{\otimes}
\newcommand{\E}{\mathop{\bf E\/}}

% bold random variables
\newcommand{\bu}{\boldsymbol{u}}
\newcommand{\bv}{\boldsymbol{v}}
\newcommand{\bx}{{\boldsymbol{x}}}
\newcommand{\by}{\boldsymbol{y}}

% T1 so literal < and > in \texttt placeholders typeset as themselves
\usepackage[T1]{fontenc}

% bra-ket
\usepackage{braket}

% --- Example structured spaces ---
\newcommand{\polyfunc}[3]{\mathcal{P}(#1, #2, #3)}
\newcommand{\polymeas}[3]{\mathrm{PolyMeas}(#1, #2, #3)}
\newcommand{\polysub}[3]{\mathrm{PolySub}(#1, #2, #3)}

% Approximation relations (paper uses \approx_\eps and \simeq_\zeta)
\newcommand{\approxeps}{\approx_{\varepsilon}}
\newcommand{\simgeq}{\simeq_{\zeta}}


\title{The Exponent Comparison in the Canonical Pauli Basis Test Bound}
\author{}
\date{2026-09-05}

\begin{document}
\maketitle

\section{At a glance}

\begin{itemize}
  \item \textbf{Difficulty.} Low and resolved.  One intermediate inequality
  in the source uses $k\leq m$ to compare the field exponent $k$ with the
  block parameter $m$.  This does not follow from the canonical parameter
  definition, but the claimed soundness bound follows by estimating the same
  term directly.
  \item \textbf{Estimated weight.} The completed replacement occupies roughly
  40--60 lines of formal proof: one reusable real-analytic lemma and its
  specialization in the canonical parameter calculation.  No work remains for
  this comparison.
  \item \textbf{Mathlib/project split.} Mathlib supplies the Taylor lower bound
  $y^n/n!\leq e^y$ and monotonicity of real powers.%
  \footnote{These are
  \leanid{Real.pow\_div\_factorial\_le\_exp} in
  \path{Mathlib/Analysis/Complex/Exponential.lean} and
  \leanid{Real.rpow\_le\_rpow\_of\_exponent\_le} in
  \path{Mathlib/Analysis/SpecialFunctions/Pow/Real.lean}.}
  The project packages these facts as an explicit polynomial--exponential
  envelope and applies it after unfolding the canonical parameters and
  soundness function.
  \item \textbf{Key mathematical inputs.} For a real exponent $s$, a rate
  $t>0$, and $x\geq1$, put
  $n=\lceil s\rceil_{+}=\max(\lceil s\rceil,0)$.  The Taylor bound at
  $y=(t\ln 2)x$ gives
  \[
    x^s2^{-tx}\leq \frac{n!}{(t\ln 2)^n}.
  \]
  Its canonical specialization gives the explicit constant used to control
  the third soundness term.
  Tracked in
  \href{https://github.com/Dengnifer/MIPStarRE-A/issues/16}{GitHub issue \#16};
  the parent blueprint task is \localissue{0002}.
\end{itemize}

\section{Key theorem forms}

Let $a\geq 1$ and $0<b<1$ be the universal constants, and let $\eps\geq0$ be a
failure-probability parameter.  The soundness function is
\[
  \delta_{\mathrm{qld}}(\eps,m,d,q)
    = a(md)^a\bigl(\eps^b+q^{-b}+2^{-bmd}\bigr).
\]
Write $\log_2$ for the base-two logarithm and $\ln$ for the natural logarithm.
For an integer $R\geq4$, let $c$ be the least even integer at least $(a+b)/b$, set
$k=c\lceil\log_2\log_2 R\rceil+1$ and $q=2^k$, and let $j$ be determined by
$2^j\leq c\lceil\log_2 R\rceil+1<2^{j+1}$. The remaining canonical parameters
are $m=2^j$ and $d=1$. In particular,
$m\geq (c/2)\log_2 R\geq\log_2 R$.
The canonical-parameter lemma asserts that there are universal constants
$a'\geq1$ and $0<b'\leq1$ such that, for every integer $R\geq4$ and
$\eps\geq0$,
\[
  \delta_{\mathrm{qld}}(\eps,m,d,q)
    \leq a'\bigl((\log_2 R)^{a'}\eps^{b'}+(\log_2 R)^{-b'}\bigr).
  \tag{\texttt{lem:delta-bound}}
\]
For the third summand, put
\[
  n=\lceil a+b\rceil_{+}=\max(\lceil a+b\rceil,0)
  \qquad\text{and}\qquad
  C=\frac{n!}{(b\ln 2)^n}.
\]
The required estimate is
\[
  a m^a 2^{-bm} \leq aC(\log_2 R)^{-b}.
\]
The constant $C$ depends only on $a$ and $b$.  This estimate, together with
the source bounds for the first two summands, gives the asserted
polylogarithmic bound for $\delta_{\mathrm{qld}}$ with $b'=b$ and
$a'=a((2c)^a+C)$.

\section{The source comparison}

The canonical parameters are defined in \cite[lines~1503--1514]{Ji2020MIPStar},
and the bound and its proof occupy lines~1520--1562.  At line~1543 the proof
uses%
\footnote{Tracked in
\href{https://github.com/Dengnifer/MIPStarRE-A/issues/16}{GitHub issue \#16}.
The canonical parameter definition is at
\path{references/qpbt-paper/08_classical_and_quantum_low_degree_tests.tex},
lines~1503--1514; the bound and its proof are at lines~1520--1562, with the
disputed comparison at line~1543.}
\[
  2^{-bmd}\leq q^{-b}
  \qquad\text{because}\qquad
  k\leq m.
  \tag{\text{source line 1543}}
\]
That comparison does not follow from the stated parameter choice.  For example,
the allowed values $a=1$ and $b=1/2$ give $c=4$, and $R=5$ gives $k=9$ but
$m=8$.  Thus the displayed justification is false even though the lemma's
conclusion remains valid.

The third term can instead be bounded independently of $q$.  More generally,
for real $s$, $t>0$, and $x\geq1$, let
$n=\lceil s\rceil_{+}=\max(\lceil s\rceil,0)$.  Since $x^s\leq x^n$, the
Taylor lower bound $y^n/n!\leq e^y$, applied at $y=(t\ln 2)x$, gives
\[
  x^s2^{-tx}
    \leq x^n2^{-tx}
    \leq \frac{n!}{(t\ln 2)^n}.
\]
Now take $s=a+b$, $t=b$, and $x=m$.  Since $a+b>0$, the constant in this
bound is the quantity $C$ above.  Moreover, $m\geq\log_2 R\geq1$ and $b>0$,
so $m^{-b}\leq(\log_2 R)^{-b}$.  Hence
\[
  a m^a2^{-bm}
    =a\bigl(m^{a+b}2^{-bm}\bigr)m^{-b}
    \leq aC(\log_2 R)^{-b}.
\]
This is the required estimate with an explicit constant; no supremum argument
is needed.

\section{Comparison with the blueprint and formalization}

The blueprint makes explicit the nonnegative failure-probability domain
implicit in the Pauli basis test soundness theorem and otherwise retains the
source statement.
The blueprint proof records the Taylor envelope above as a separate
formalization-support lemma and uses the same explicit constant $C$.

The formal development proves the general envelope and specializes it in the
canonical parameter calculation.  The canonical soundness bound is now fully
proved, so no formal proof obligation remains; its additional nonnegativity
premise is only the faithful domain condition on the failure parameter.%
\footnote{The envelope is
\leanid{MIPStarRE.LDT.rpow\_mul\_two\_rpow\_neg\_le} in
\doclink{MIPStarRE/LDT/Basic/RpowBounds.lean}; its specialization proves
\leanid{MIPStarRE.QPBT.exists\_deltaQld\_introParams\_bound} in
\doclink{MIPStarRE/QPBT/Test/CanonicalParams.lean}.}

\section{Conclusion}

The source's comparison $k\leq m$ is false for some permitted constants and
values of $R$.  This is a local defect in the proof, not a gap in the lemma's
statement: the Taylor envelope bounds $m^{a+b}2^{-bm}$ uniformly, and the
canonical lower bound $m\geq\log_2 R$ supplies the remaining negative power.
This replacement is now formalized and proves the source conclusion after
enlarging the universal constant.

\bibliographystyle{alpha}
\bibliography{references}

\end{document}
